\section{为什么数据最重要}

本节先回答整讲的动机：当 architecture、optimizer 与并行 recipe 越来越公开时，数据来源、配比和过滤规则反而成为最难观察的变量。模型读过什么决定知识覆盖，哪些样本被删除决定语言与价值偏好，provenance 是否完整则决定结果能否审计和持续使用。

\begin{figure}[H]
\centering
\includegraphics[width=0.92\textwidth]{llama3-data.png}
\caption{Llama 3 报告中架构和训练流程相对透明，但数据细节非常有限。}
\figsource{CS336 Lecture 13 官方源引用图。}
\end{figure}

\begin{knowledgebox}{讲义提醒：看到模型差异时，先别只归因于架构}
两个参数量和训练 FLOPs 相近的模型仍可能因为数据语言比例、代码/数学含量、重复率和质量阈值不同而表现迥异。没有 data card 与可复现 pipeline 时，很多“架构优势”只是无法排除数据混杂变量。这是根据课程“数据是模型差异化核心变量”的主线进一步得到的审计建议，而不是对某个具体模型的因果断言。
\end{knowledgebox}

\begin{knowledgebox}{读图：为什么模型公司对 data 更保密}
开源权重模型往往公开 architecture、parameter count、training recipe，甚至并行策略；但数据来源、比例、清洗规则、过滤器和法律授权常只给模糊描述。原因包括竞争优势和版权责任。数据成为模型差异化的核心资产。
\end{knowledgebox}

\begin{knowledgebox}{术语消化：训练阶段和数据质量}
\begin{tabular}{L{0.22\textwidth}L{0.36\textwidth}L{0.32\textwidth}}
\toprule
阶段 & 数据形态 & 目标 \\
\midrule
Pre-training & 大量原始或轻清洗文本、代码、网页、书籍等。 & 建立通用语言和世界知识。 \\
Mid-training & 更高质量、能力导向的数据。 & 强化代码、数学、推理、多语等能力。 \\
Post-training & 聊天、偏好、RL、工具使用轨迹。 & 对齐交互行为和任务偏好。 \\
\bottomrule
\end{tabular}
\end{knowledgebox}

\begin{figure}[H]
\centering
\includegraphics[width=0.78\textwidth]{olmo2-pretraining.png}
\caption{OLMo 2 pre-training 数据示例。}
\figsource{CS336 Lecture 13 官方源引用图。}
\end{figure}

\begin{figure}[H]
\centering
\includegraphics[width=0.7\textwidth]{olmo2-dolmino.png}
\caption{OLMo 2 mid-training / Dolmino 数据示例。}
\figsource{CS336 Lecture 13 官方源引用图。}
\end{figure}

\begin{figure}[H]
\centering
\includegraphics[width=0.72\textwidth]{tulu.png}
\caption{Tulu post-training 数据示例。}
\figsource{CS336 Lecture 13 官方源引用图。}
\end{figure}

\begin{knowledgebox}{读图：OLMo/Tulu 三张图的教学含义}
三张图展示了从大量弱筛网页和通用语料，到更有针对性的高质量数据，再到 instruction/chat 数据的移动。数据量通常下降，人工或模型筛选强度上升，目标也从“预测文本”转向“表现出想要的行为”。
\end{knowledgebox}

